% !TEX root = ../my-thesis.tex
%
\chapter{Prompts, Reproducibility, and Proof Details}
\label{sec:appendix}

\section{Prompts}
\label{sec:appendix:prompts}

This appendix lists the prompts of the final configuration verbatim, in
the versions used by the reference evaluation run of
Section~\ref{sec:evaluation:gold} (\texttt{classify\_batch\_v4} and
\texttt{verify\_v4}). Placeholders in curly braces are filled by the pipeline
at run time. The typed classification prompt and the verifier prompt contain
the two rebuttal rules discussed in Section~\ref{sec:method:extraction} as
well as the general not-an-attack rule of the final revision. The cross-topic
and ablation batch of Sections~\ref{sec:evaluation:crosstopic}
and~\ref{sec:evaluation:ablation} used the preceding version
(\texttt{classify\_batch\_v3} and \texttt{verify\_v3}), which states the
stance rule only inside the rebuttal definition and lacks the general
not-an-attack test and the specificity requirements for undercutting and
undermining. The binary prompt (\texttt{classify\_batch\_binary\_v1}) is
identical in both batches and omits the
attack type, the type definitions, and these rules.

\subsection*{Atomization prompt}

\begin{lstlisting}[breaklines=true, basicstyle=\ttfamily\footnotesize]
You receive an argument from a debate on "{topic}".
Check whether it bundles multiple distinct claims into one sentence.
If yes, split it into atomic arguments (one claim each), preserving the stance.
If it already expresses a single claim, return it unchanged.

Important: Only promote claims the speaker actually ASSERTS as their own.
- Concessive/contrastive clauses (markers: "but", "although", "while",
  "despite", "even though", "yet") are rhetorical hedges, NOT independent
  arguments. Keep them attached to the main claim or drop them - never make
  them a separate atom.
- Reported or conceded positions of the opposing side (e.g. "opponents argue
  X", "greens always argue X", "some claim X, but ...") are NOT the speaker's
  own claim. Do NOT promote them to a separate atom; drop them.
This prevents atoms that, read in isolation, carry the OPPOSITE stance of the
speaker (e.g. a pro-nuclear sentence must never yield a pro-renewable atom).

Self-containment rule: Every output argument MUST be understandable without
any surrounding context. If a fragment contains pronouns like "it", "them",
"this", "they" without a clear referent, replace the pronoun with the explicit
noun from the original sentence. If the referent cannot be recovered, drop
the fragment entirely.

Argument [{id}] ({stance}):
{text}

Respond ONLY with a JSON array, no markdown. Use IDs like "{id}", "{id}b", ...:
[
  {"id": "{id}",  "text": "...", "stance": "{stance}"},
  {"id": "{id}b", "text": "...", "stance": "{stance}"}
]
\end{lstlisting}

\subsection*{Typed classification prompt (batched)}

\begin{lstlisting}[breaklines=true, basicstyle=\ttfamily\footnotesize]
Propose attack candidates for the following argument pairs from a debate on
"{topic}".

An attack exists when one argument:
- asserts a conclusion logically incompatible with the other's conclusion -
  the one claim can be true only if the other is false (rebuttal), OR
- undermines a premise the other relies on (undermining), OR
- shows the other's conclusion does not follow even if its premise is true
  (undercutting)

NOT an attack - this applies to ALL three types: merely taking the opposing
stance, addressing a different sub-topic, or generally "arguing for the other
side". A con argument does not attack every pro argument (and vice versa).
Test: if both specific claims can be true at the same time without tension,
there is NO attack and direction MUST be NONE.

{pairs}

direction    = exactly one of: A_ATTACKS_B / B_ATTACKS_A / NONE
attack_type  = one of: rebuttal / undercutting / undermining / NONE
  rebuttal:     the two conclusions are logically incompatible - they negate
                the SAME proposition, so one claim is true only if the other
                is false; there must be a direct contradiction between the two
                specific claims. A descriptive/empirical claim (what people
                believe or do, statistics, trends) and a normative claim (what
                ought to be) are NOT logically incompatible - they cannot
                rebut each other (at most undermining/undercutting). A
                rebuttal is inherently MUTUAL (it holds in both directions);
                pick either direction.
  undercutting: attacker shows the target's conclusion doesn't follow even if
                the target's premise holds - it targets the INFERENCE of that
                specific argument. A general counter-consideration on the same
                topic is NOT undercutting.
  undermining:  attacker directly attacks a premise the target relies on. The
                premise must be a SPECIFIC claim the target states or
                presupposes - the target's overall stance or the general
                attractiveness of its position is NOT a premise. In "reason",
                name the attacked premise; if you cannot name it, the type is
                not undermining.
  If direction=NONE, attack_type MUST be NONE.
confidence   = integer 0-100: probability that ANY attack exists between the
  two arguments (0 = definitely no attack, 100 = definitely attack)
  If direction=NONE, confidence MUST be 0.

Respond ONLY with a JSON array, same order as the pairs above, no markdown:
[
  {"pair": "id_a:id_b", "direction": "A_ATTACKS_B", "attack_type": "rebuttal",
   "confidence": 85, "reason": "one sentence"},
  ...
]
\end{lstlisting}

\subsection*{Binary classification prompt (batched, ablation arm)}

\begin{lstlisting}[breaklines=true, basicstyle=\ttfamily\footnotesize]
Propose attack candidates for the following argument pairs from a debate on
"{topic}".

An attack exists when one argument is logically incompatible with the other's
conclusion, undermines a premise the other relies on, or shows that the
other's conclusion does not follow even if its premise holds. Merely taking
the opposing stance, addressing a different sub-topic, or arguing for the
other side is NOT enough. There must be a genuine conflict between the two
specific claims.

{pairs}

direction    = exactly one of: A_ATTACKS_B / B_ATTACKS_A / NONE
confidence   = integer 0-100: probability that ANY attack exists between the
  two arguments (0 = definitely no attack, 100 = definitely attack)
  If direction=NONE, confidence MUST be 0.

Respond ONLY with a JSON array, same order as the pairs above, no markdown:
[
  {"pair": "id_a:id_b", "direction": "A_ATTACKS_B", "confidence": 85,
   "reason": "one sentence"},
  ...
]
\end{lstlisting}

\subsection*{Verifier prompt (typed)}

\begin{lstlisting}[breaklines=true, basicstyle=\ttfamily\footnotesize]
Two arguments from a debate on "{topic}":

[{id_a}] ({stance_a}) {text_a}
[{id_b}] ({stance_b}) {text_b}

A previous analysis claims: {claimed}

Does this attack relation genuinely hold?
Does A undermine a premise B relies on (undermining)? The premise must be a
SPECIFIC claim B states or presupposes - B's overall stance or the general
attractiveness of its position is NOT a premise. If you cannot name the
attacked premise, it is not undermining.
Do A and B assert logically incompatible conclusions that negate the same
proposition, so that one is true only if the other is false (rebuttal)?
Opposing stance alone is NOT enough.
A descriptive claim (what people believe/do, statistics) and a normative claim
(what ought to be) are NOT logically incompatible and cannot rebut each other.
Does A show B's conclusion doesn't follow even if B's premise is true
(undercutting)?
A general counter-consideration on the same topic is NOT undercutting.
Decisive test: if both specific claims can be true at the same time without
tension, there is NO attack - merely being on opposite sides of the debate is
not enough.
If any of the three genuinely holds, confirm=true.

Respond ONLY with JSON, no markdown:
{"confirms": true, "reason": "one sentence"}
\end{lstlisting}

\section{Reproducibility}
\label{sec:appendix:repro}

All experiments use the following canonical configuration. The model is
GPT-OSS-120B~\cite{gptoss2025}, called through the university-hosted inference
endpoint at temperature zero with a batch size of ten pairs per classification
call. Per
topic, 14 arguments are sampled from the UKP test split, seven pro and seven
con, with random seed 42. Candidates are accepted at confidence 70 or above,
verified between 50 and 69, and rejected below 50. Rebuttal symmetrization is
enabled in the typed arm and structurally impossible in the binary arm. The
solver computes preferred and stable extensions with the grounded-first
reduction and skips them if the undecided core exceeds 18 arguments, which did
not occur in the evaluation.

Every run writes a log directory containing the run configuration, the frozen
atomized argument set, one record per queried pair with the model output and
routing outcome, the accepted attacks, the generated Prolog facts, and the
solver results. The ablation additionally stores, per topic, the shared atom
set and one full run log per extraction arm, together with a comparison file
holding the edge sets and Jaccard values. All tables in
Chapter~\ref{sec:evaluation} are generated by a single script from these
artifacts, so that every reported number is reproducible from the logged runs
without further API calls.

The complete implementation, the run logs of all reported experiments, the
reference annotation, and the table generation script are publicly available
in the accompanying repository.\footnote{\url{https://github.com/AaronHabyarimana/structuring-llm-arguments}}
The UKP corpus as a whole is not redistributed and must be
obtained from its original authors. The run logs contain only the small
per topic samples of sentences that were drawn for the experiments, as part
of the experimental record. The stance annotations included in these logs
are released by the corpus authors under a Creative Commons Attribution
NonCommercial (CC BY-NC) license~\cite{stab2018}.

\section{Algorithm and Proof Sketch for the Grounded-First Reduction}
\label{sec:appendix:proof}

Algorithm~\ref{alg:method:solver} summarises the grounded-first solving
procedure of Section~\ref{sec:method:solver}.

\begin{algorithm}[tb]
\caption{Grounded-first solving with undecided-core guard}
\label{alg:method:solver}
\begin{algorithmic}[1]
\Require framework $AF = \langle \mathcal{A}, \mathcal{R} \rangle$,
  bound $\mathit{MaxU}$
\Ensure grounded extension $G$, preferred set $\mathcal{P}$,
  stable set $\mathcal{S}$
\State $G \gets \emptyset$
\Repeat \Comment{fixed-point iteration, polynomial}
  \State $G \gets F_{AF}(G)$
\Until{$G$ unchanged}
\State $U \gets \{x \in \mathcal{A} \mid x \notin G \text{ and no }
  g \in G \text{ attacks } x\}$
\If{$|U| > \mathit{MaxU}$}
  \State \Return $(G, \bot, \bot)$
    \Comment{enumeration skipped, only grounded reported}
\EndIf
\State $\mathcal{P}_U \gets$ maximal admissible subsets of $AF|_U$, with
  attacks from outside $U$ discounted as neutralised
\State $\mathcal{S}_U \gets$ conflict-free $E \subseteq U$ that attack every
  $x \in U \setminus E$
\State \Return $(G,\; \{G \cup E \mid E \in \mathcal{P}_U\},\;
  \{G \cup E \mid E \in \mathcal{S}_U\})$
\end{algorithmic}
\end{algorithm}

The rest of this section gives the proof sketch for
Proposition~\ref{prop:method:reduction}, which states that the
$\sigma$-extensions of $AF$ for
$\sigma \in \{\text{complete}, \text{preferred}, \text{stable}\}$ are exactly
the sets $G \cup E$ where $E$ is a $\sigma$-extension of the framework
restricted to the undecided core $U$.

\begin{proof}[Proof sketch]
Let $S$ be a complete extension of $AF$. Since the grounded extension is the
least complete extension, $G \subseteq S$, and conflict-freeness excludes from
$S$ every argument attacked by $G$, so $S = G \cup E$ for some $E \subseteq U$.
By definition of $U$, no attacker of an argument in $U$ lies in $G$, so every
attacker of an argument in $E$ is either attacked by $G$, and hence neutralised
in $S$, or lies in $U$ itself. Defense within $S$ therefore coincides with
defense within the discounted restriction to $U$. In the other direction,
$G \cup E$ is conflict-free for every conflict-free $E \subseteq U$, because an
attacker of a member of $G$ is itself attacked by $G$ and therefore lies
outside $U$, and the same case analysis shows that $G \cup E$ is complete in
$AF$ whenever $E$ is complete in the restriction. Together,
$S \mapsto S \setminus G$
is a one-to-one correspondence between the complete extensions of $AF$ and
those of the restriction. Preferred extensions are the maximal complete ones,
and the correspondence preserves inclusion, which settles the preferred case.
For the stable case, an argument outside $G \cup E$ is either attacked by $G$
or lies in $U \setminus E$, where stability of $E$ in the restriction requires
an attack from $E$, so $G \cup E$ attacks every outside argument, and the
converse direction follows by restricting a stable extension of $AF$ to $U$.
\end{proof}
